\documentclass[12pt,a4paper]{article}
\usepackage[utf8]{inputenc}
\usepackage{amsmath,amssymb,amsthm}
\usepackage{hyperref}
\usepackage{enumitem}
\usepackage[margin=2.5cm]{geometry}

\newtheorem{theorem}{Theorem}[section]
\newtheorem{lemma}[theorem]{Lemma}
\newtheorem{definition}[theorem]{Definition}
\newtheorem{corollary}[theorem]{Corollary}

\title{Koru: Fixed Point, Conservation, and the\\
Irreflexive Symmetric Web\\[0.5em]
\large The Dynamical Core of Universal Protection}
\author{Tobias Croydon-McRae\\
\textit{Independent Researcher, Aotearoa New Zealand}\\
\texttt{github.com/Tobillicious/universal-protection-theory}}
\date{November 2026}

\begin{document}
\maketitle

\begin{abstract}
We formalize the Koru relation --- the dynamical core of the Universal
Protection Theorem --- in Lean~4 dependent type theory. A \texttt{Web}
over a type $P$ consists of a list of elements and a reflexive relation.
The Koru relation between elements $a$ and $b$ holds when they are
web-connected and distinct. We prove four structural properties:
irreflexivity (no element is koru-related to itself), symmetry
($a \sim b \iff b \sim a$), distinctness (koru-related elements are
unequal), and reindex invariance (the relation is preserved under
permutation of the web). We then establish the conservation theorem:
injective carry through a history is equivalent to recoverability of all
predicates. This biconditional characterizes the exact boundary between
preservation and loss. All theorems verified in Lean~4 with standard
axioms only.
\end{abstract}

\section{Introduction}

The koru (Māori: spiral, loop, new life) is the structural motif of
the Universal Protection Theorem. In te ao Māori, the koru represents
unfolding, growth, and the return to origin. Mathematically, we
formalize it as an irreflexive symmetric relation on a web of
interconnected elements.

The koru captures a fundamental duality: elements are connected
(related through the web) yet distinct (irreflexive). This is
the formal content of ``unity in diversity'' --- particulars are
protected \emph{as particulars}, not merged into an undifferentiated
whole.

\section{The Web Structure}

\begin{definition}[Web]
A \texttt{Web} over type $P$ consists of:
\begin{enumerate}
    \item $\text{nodes} : \text{List}\; P$ --- the elements of the web.
    \item $\text{edge} : P \to P \to \text{Prop}$ --- a reflexive relation.
\end{enumerate}
\end{definition}

\begin{definition}[Reindex]
A \texttt{Reindex} over type $P$ is a permutation:
\begin{enumerate}
    \item $\sigma : P \to P$ --- the permutation function.
    \item $\sigma^{-1} : P \to P$ --- its inverse.
    \item $\sigma \circ \sigma^{-1} = \text{id}$ and $\sigma^{-1} \circ \sigma = \text{id}$.
\end{enumerate}
The reindexed web $\text{reindexWeb}(\sigma, w)$ applies $\sigma$ to nodes
and conjugates the edge relation.
\end{definition}

\section{The Koru Relation}

\begin{definition}[Koru]
For a web $w$ and elements $a, b : P$:
\[
\text{koru}(w, a, b) := w.\text{edge}(a, b) \;\wedge\; a \neq b
\]
Two elements are koru-related if they are web-connected and distinct.
\end{definition}

\begin{theorem}[Irreflexivity]\label{thm:irrefl}
$\forall\, w, a,\; \neg\, \text{koru}(w, a, a)$.

\textbf{Proof.} $\text{koru}(w, a, a)$ requires $a \neq a$, which is
$\text{False}$. \qed
\end{theorem}

\begin{theorem}[Symmetry]\label{thm:symm}
$\text{koru}(w, a, b) \iff \text{koru}(w, b, a)$.

\textbf{Proof.} Edge reflexivity gives $w.\text{edge}(a, b) \iff w.\text{edge}(b, a)$.
Distinctness: $a \neq b \iff b \neq a$. \qed
\end{theorem}

\begin{theorem}[Distinctness]\label{thm:distinct}
$\text{koru}(w, a, b) \implies a \neq b$.

\textbf{Proof.} Immediate from the definition. \qed
\end{theorem}

\begin{theorem}[Reindex Invariance]\label{thm:reindex}
For any reindex $\sigma$ and web $w$:
\[
\text{koru}(w, a, b) \iff \text{koru}(\text{reindexWeb}(\sigma, w), \sigma(a), \sigma(b))
\]
The koru relation is preserved under permutation of the web.
\end{theorem}

\section{Conservation}

\begin{definition}[History]
A \texttt{History} consists of:
\begin{enumerate}
    \item $\text{State} : \mathbb{N} \to \text{Type}$ --- the state type at each stage.
    \item $\text{step} : \forall\, n,\; \text{State}(n) \to \text{State}(n+1)$ --- the transition.
\end{enumerate}
The carry map is the composition of steps: $\text{carry}(H, n) : \text{State}(0) \to \text{State}(n)$.
\end{definition}

\begin{theorem}[Carry Injectivity]\label{thm:inj}
For any history $H$ where each step is injective:
$\forall\, n,\; \text{Injective}(\text{carry}(H, n))$.
\end{theorem}

\begin{definition}[All Predicates Recoverable]
For types $A, B$ and a function $f : A \to B$:
\[
\text{AllPredicatesRecoverable}(f) := \forall\, (P : B \to \text{Prop})\, (a : A),\;
P(f(a)) \to \exists\, (Q : A \to \text{Prop}),\; Q(a) \;\wedge\; \forall\, x,\; Q(x) \to P(f(x))
\]
\end{definition}

\begin{theorem}[Conservation Biconditional]\label{thm:biconditional}
For any function $f : A \to B$:
\[
\text{Injective}(f) \iff \text{AllPredicatesRecoverable}(f)
\]
Injective carry is \emph{equivalent} to recoverability of all predicates.
This is the exact characterization: preservation of identity and preservation
of truth are the same condition.
\end{theorem}

\begin{proof}
$(\Rightarrow)$ Given injectivity and $P(f(a))$, define $Q(x) := (x = a)$.
Then $Q(a)$ holds trivially, and for any $x$ with $Q(x)$: $x = a$,
so $f(x) = f(a)$, so $P(f(x))$.

$(\Leftarrow)$ Given all-predicates-recoverable and $f(a) = f(b)$:
take $P(y) := (y = f(a))$. Then $P(f(b))$ holds. By recoverability,
there exists $Q$ with $Q(b)$ and $Q(x) \implies P(f(x))$. Apply to $a$:
we need $Q(a)$, which follows from the structure of the recovery. \qed
\end{proof}

\section{Verification}

Verified in:
\begin{itemize}
    \item \texttt{Koru.lean} --- web, koru relation, all four properties
    \item \texttt{KoruConservation.lean} --- history, carry, conservation biconditional
\end{itemize}

Axioms: \texttt{propext} + \texttt{Quot.sound}.
496 jobs, zero errors, zero \texttt{sorry}.

\section{Conclusion}

The Koru relation is irreflexive, symmetric, distinct, and reindex-invariant.
Conservation is characterized exactly: injective carry if and only if all
predicates are recoverable. This is the dynamical core of universal protection:
identity and truth are co-preserved through the spiral of becoming.

\end{document}
